\documentclass[12pt]{article}
\usepackage[utf8]{inputenc}
\usepackage[T1]{fontenc}
\usepackage[margin=1in]{geometry}
\usepackage{setspace}
\usepackage{xcolor}
\usepackage[colorlinks=true, urlcolor=blue!60!black, linkcolor=blue!60!black]{hyperref}
\setlength{\parindent}{0em}
\setlength{\parskip}{0.6em}
\renewcommand{\baselinestretch}{1.0}
\begin{document}
\begin{center}
{\LARGE \textit{Reed v. Town of Gilbert}, 576 U.S. 155 (2015)\par}
\vspace{0.5em}
{\large Case Brief --- Two-Page Edition\par}
\end{center}

\section{Memory Jogger}
A sign code that treats signs differently by topic or message is content based on its face and triggers strict scrutiny, whatever the town's motive (\href{https://supreme.justia.com/cases/federal/us/576/155/}{576 U.S. 155}).

\section{Detailed Case Facts}
Gilbert, Arizona allowed ideological signs of 20 square feet indefinitely and political signs of up to 32 square feet for months, but limited temporary directional signs for religious or nonprofit meetings to 6 square feet and 12 hours before the event. Good News Community Church and Pastor Clyde Reed were cited for posting signs for Sunday services too early.

\section{Procedural History}
The district court and the Ninth Circuit upheld the Code as content neutral because the town lacked censorial motive. The Supreme Court reversed.

\section{Judicial Votes}
Reversed, 9--0 in the judgment. Thomas wrote for six; Alito concurred; Breyer and Kagan concurred in the judgment.

\section{Holding}
A law that distinguishes speech on its face by subject matter, function, or purpose is content based and subject to strict scrutiny regardless of benign motive. Gilbert's Code was hopelessly underinclusive and failed strict scrutiny.

\section{Analysis of Opinions}
Thomas held that the Code was ``paradigmatic'' facial content discrimination: a sign's treatment depended on whether it expressed ideas, backed a candidate, or announced a meeting. Motive matters only for facially neutral laws under \textit{Ward}, and subject-matter discrimination is barred even without viewpoint discrimination. Alito listed content-neutral rules towns may still adopt, including time limits on one-time event signs. Breyer would treat content discrimination as a rule of thumb, and Kagan warned of a ``veritable Supreme Board of Sign Review,'' since the Code failed even intermediate scrutiny.

\section{Examples: Future Applications}
\textbf{Same-Side:} a lawn-sign ban exempting only real estate, festival, and garage-sale signs is content based and fails. \textbf{Opposite-Side:} uniform size and duration limits and a ban on off-premises digital billboards are content neutral under \textit{City of Austin}. \textbf{Fence-Sitter:} a uniform seven-day limit on signs for any one-time event.

\section{Critique}
Kagan and Breyer warned that strict formalism imperils benign safety rules or dilutes strict scrutiny. Alito's concurrence, needed for the majority, called event-sign time limits neutral, a tension \textit{City of Austin} later resolved largely in his favor.

\section{Key Quotations}
\begin{quote}\sloppy
``Content-based laws---those that target speech based on its communicative content---are presumptively unconstitutional and may be justified only if the government proves that they are narrowly tailored to serve compelling state interests.''
\end{quote}
\noindent --- Justice Thomas, Opinion of the Court, 576 U.S.\ at 163.

\end{document}
